% year: תשפ"ב
% type: מבחן
% semester: א
% moed: מיוחד
% number: 6א
% points: 8
% categories: integral-applications
% source: exams/מבחנים/תשפב/מועד ג תשפב.pdf

\begin{question}
(8 נק') חשבו את השטח הסופי החסום בין גרף הפונקציה $f(x)=x^3-x$ וציר ה-$x$.
\end{question}

\begin{hint}
שימו לב שהפונקציה מחליפה סימן בין נקודות החיתוך עם הציר — יש לחשב את השטח של כל חלק בנפרד (או להשתמש באי-זוגיות).
\end{hint}

\begin{solution}
\textbf{חיתוך עם ציר ה-$x$:} $x^3-x=x(x-1)(x+1)=0\iff x\in\{-1,0,1\}$.
בקטע $(-1,0)$ מתקיים $f>0$, ובקטע $(0,1)$ מתקיים $f<0$. מחוץ ל-$[-1,1]$ התחום בין הגרף לציר אינו חסום, ולכן השטח הסופי החסום הוא שני ה"עלים" שבין $-1$ ל-$1$:
\[ S=\int_{-1}^0(x^3-x)\,dx+\int_0^1(x-x^3)\,dx. \]
\[ \int_0^1(x-x^3)\,dx=\left[\frac{x^2}{2}-\frac{x^4}{4}\right]_0^1=\frac14,\qquad \int_{-1}^0(x^3-x)\,dx=\left[\frac{x^4}{4}-\frac{x^2}{2}\right]_{-1}^0=-\left(\frac14-\frac12\right)=\frac14. \]
(השוויון נובע גם מכך ש-$f$ אי-זוגית.) לכן
\[ S=\frac14+\frac14=\boxed{\frac12}. \]
(שימו לב: $\int_{-1}^1(x^3-x)\,dx=0$ — זה אינו השטח.)
\end{solution}
